\answer{ex:22:1}{(a)~$32\%$ y $100\%$. (b)~$R_\uparrow-R_\downarrow\ge2\sqrt{2000\,\text{Hz}/1\,\text{s}}\approx89$~Hz, apenas el $4.5\%$ de la frecuencia total.}
\answer{ex:22:2}{(a)~$n\ge4(p_1q_1+p_2q_2)/\Delta^2\approx560$. (b)~$\approx56$.}
\answer{ex:22:3}{(a)~Balance detallado: $\rho PK$ es simétrico. (b)~La fracción de fallos es la media armónica $1/\langle1/P\rangle=0.18$, frente a $0.5$ sin retroalimentación.}
\answer{ex:22:4}{(a)~Un paralelogramo de área $4h^2=0.04$ (teniendo en cuenta la limitación de $p$, numéricamente $0.045$). (b)~$\approx0.18$.}
\answer{ex:22:5}{$0.49$: parte de los cultivos se van a una región donde casi siempre fallan y vagan por ella. Con la limitación, $1.00$: en un conjunto acotado la densidad $\propto1/P$ se concentra en la región absorbente.}
\answer{ex:22:6}{(a)~Hacen falta $\ge5$ niveles; $8$ electrodos bastan. (b)~No: haría falta $r_{\max}\ge25/T=100$~Hz.}
\answer{ex:22:7}{Problema abierto. El tiempo de la búsqueda aleatoria crece con $d$ aproximadamente como la razón de volúmenes $(L/h)^d$; el de la búsqueda con signo del error, linealmente en $d$. Separa las hipótesis un experimento con permutación: tras el fallo, un estímulo predecible pero fuerte; tras el acierto, uno impredecible pero débil; hacen falta varias decenas de cultivos en cada grupo.}
